\documentclass[12pt]{article}
\usepackage{amsmath, amssymb}
\usepackage{geometry}
\usepackage{setspace}
\usepackage{enumitem}
\geometry{margin=1in}

\title{Lecture Scribe \\ Joint Random Variables and Joint Distributions}
\date{}

\begin{document}

\maketitle
\onehalfspacing

\section{Motivation: Why Study Joint Random Variables?}

In many scientific and real-world problems, we are interested in two or more random variables simultaneously rather than analyzing them individually. Studying them together allows us to understand the relationship or dependence between the variables.

\subsection*{Examples of Paired Random Variables}

Typical examples include:

\begin{itemize}
\item Height and weight of giraffes
\item IQ and birthweight of children
\item Frequency of exercise and rate of heart disease
\item Air pollution level and respiratory illness rate
\item Number of Facebook friends and age of members
\end{itemize}

These examples illustrate situations where two variables interact or influence each other.

\subsection*{Engineering and Computing Applications}

Several practical applications illustrate the importance of studying joint random variables.

\subsubsection*{Example 1: Smart Transportation}

Two random variables are considered:

\[
X = \text{Vehicle Speed}
\]

\[
Y = \text{Traffic Density}
\]

Key question:

Can vehicle speed be high when traffic density is also high?

This analysis helps in:

\begin{itemize}
\item Adaptive traffic signal systems
\item Congestion prediction
\item Autonomous vehicle routing
\end{itemize}

Because the values of $X$ and $Y$ influence each other, they are correlated random variables.

\subsubsection*{Example 2: Wireless Network Performance}

Two important variables are:

\[
X = \text{Signal-to-Noise Ratio (SNR)}
\]

\[
Y = \text{Data Throughput}
\]

Throughput often depends on SNR. However, even with high SNR, throughput may not always be high due to network congestion. This illustrates conditional relationships between variables.

Applications include:

\begin{itemize}
\item 5G scheduling
\item Adaptive modulation
\item Network planning
\end{itemize}

\subsection*{Additional Examples}

Weather prediction:

\[
X = \text{Rainfall}, \quad Y = \text{Temperature}
\]

Drone delivery:

\[
X = \text{Wind speed}, \quad Y = \text{Delivery time}
\]

Dice experiment:

Two dice outcomes analyzed jointly

These examples demonstrate that many real-world problems require analyzing multiple random variables simultaneously.

\subsection*{Role of Joint Distributions}

When random variables appear together, they follow a joint distribution.

The joint distribution allows us to:

\begin{itemize}
\item Compute probabilities involving multiple variables
\item Understand relationships between variables
\end{itemize}

If the variables are independent, analysis becomes simpler. Otherwise, dependence is analyzed using measures such as:

\begin{itemize}
\item Covariance
\item Correlation
\end{itemize}

\section{Discrete Case: Joint Probability Mass Function (Joint PMF)}

\subsection*{Setup}

Suppose two discrete random variables:

\[
X = \{x_1, x_2, ..., x_n\}
\]

\[
Y = \{y_1, y_2, ..., y_m\}
\]

The ordered pair $(X,Y)$ takes values from the product set:

\[
\{(x_1,y_1),(x_1,y_2),...,(x_n,y_m)\}
\]

Each element represents a joint outcome.

\subsection*{Definition: Joint PMF}

The joint probability mass function is defined as:

\[
p(x_i,y_j) = P(X=x_i,Y=y_j)
\]

This represents the probability that:

\begin{itemize}
\item $X$ equals $x_i$, and
\item $Y$ equals $y_j$
\end{itemize}

simultaneously.

\subsection*{Joint PMF Table Representation}

The joint PMF is often organized as a table where:

\begin{itemize}
\item Rows correspond to values of $X$
\item Columns correspond to values of $Y$
\end{itemize}

Each entry contains:

\[
p(x_i,y_j)
\]

\subsection*{Key Properties}

Probability values lie between 0 and 1

\[
0 \le p(x_i,y_j) \le 1
\]

Total probability equals 1

\[
\sum_{i=1}^{n} \sum_{j=1}^{m} p(x_i,y_j) = 1
\]

These conditions ensure the function is a valid probability distribution.

\section{Example: Rolling Two Dice}

\subsection*{Experiment}

Two fair dice are rolled.

Define:

\[
X = \text{value on the first die}
\]

\[
T = \text{total (sum) of both dice}
\]

Each pair $(i,j)$ occurs with probability:

\[
\frac{1}{36}
\]

since there are 36 equally likely outcomes.

\subsection*{Joint Probability Table}

The joint table for $(X,T)$ lists probabilities for combinations such as:

\[
(X=1,T=2),(X=1,T=3),...,(X=6,T=12)
\]

Entries in the table are either:

\[
0 \quad \text{or} \quad \frac{1}{36}
\]

depending on whether the combination is possible.

\subsection*{Observation}

The variables $X$ and $T$ are not independent.

Reason:

The value of $T$ depends on the value of $X$.

If $X=6$, then $T$ cannot be less than 7.

Thus the variables are dependent random variables.

\section{Continuous Case: Joint Probability Density Function (Joint PDF)}

The continuous case is analogous to the discrete case, with the following replacements:

\begin{center}
\begin{tabular}{|c|c|}
\hline
Discrete & Continuous \\
\hline
Discrete values & Continuous intervals \\
Joint PMF & Joint PDF \\
Summation & Double integral \\
\hline
\end{tabular}
\end{center}

\subsection*{Product Region}

If:

\[
X \in [a,b], \quad Y \in [c,d]
\]

then the pair $(X,Y)$ lies in the product region

\[
[a,b] \times [c,d]
\]

This region represents all possible joint values.

\subsection*{Definition: Joint PDF}

A function $f(x,y)$ is the joint probability density function if:

\[
P((X,Y) \text{ lies in a small rectangle}) \approx f(x,y)\,dx\,dy
\]

where the rectangle has area:

\[
dx \times dy
\]

Thus:

\[
\text{Probability} \approx f(x,y)dxdy
\]

This is illustrated geometrically as a small rectangle in the $(x,y)$ plane.

\subsection*{Key Properties of Joint PDF}

Non-negativity

\[
f(x,y) \ge 0
\]

Total probability equals 1

\[
\int_{c}^{d}\int_{a}^{b} f(x,y)\,dx\,dy = 1
\]

Important note:

The value of $f(x,y)$ may exceed 1 because it is a density, not a probability.

\section{Worked Example (Continuous Joint PDF)}

\subsection*{Given}

\[
f(x,y)=4xy, \quad 0 \le x \le 1, \quad 0 \le y \le 1
\]

Thus $(X,Y)$ lies in the unit square

\[
[0,1]^2
\]

Tasks:

\begin{itemize}
\item Verify that $f(x,y)$ is a valid joint PDF
\item Visualize the event
\[
A=\{X<0.5,Y>0.5\}
\]
\item Compute $P(A)$
\end{itemize}

\subsection*{Step 1: Check Validity (Normalization)}

Compute total probability:

\[
\int_{0}^{1}\int_{0}^{1} 4xy\,dx\,dy
\]

First integrate with respect to $x$:

\[
\int_{0}^{1}4xy\,dx = 4y\left[\frac{x^2}{2}\right]_{0}^{1}
\]

\[
=2y
\]

Now integrate with respect to $y$:

\[
\int_{0}^{1}2y\,dy
\]

\[
=[y^2]_{0}^{1}
\]

\[
=1
\]

Since:

\[
f(x,y) \ge 0
\]

total probability = 1

the function is a valid joint PDF.

\subsection*{Step 2: Compute $P(A)$}

Event:

\[
A=\{X<0.5,Y>0.5\}
\]

Thus:

\[
P(A)=\int_{0}^{0.5}\int_{0.5}^{1}4xy\,dy\,dx
\]

First integrate with respect to $y$:

\[
\int_{0.5}^{1}4xy\,dy
\]

\[
=4x\left[\frac{y^2}{2}\right]_{0.5}^{1}
\]

\[
=2x(1-0.25)
\]

\[
=\frac{3}{2}x
\]

Now integrate with respect to $x$:

\[
\int_{0}^{0.5}\frac{3}{2}x\,dx
\]

\[
=\frac{3}{2}\left[\frac{x^2}{2}\right]_{0}^{0.5}
\]

\[
=\frac{3}{16}
\]

Therefore:

\[
P(A)=\frac{3}{16}
\]

\section{Joint Cumulative Distribution Function (Joint CDF)}

Suppose $X$ and $Y$ are jointly distributed random variables.

The notation:

\[
X \le x, \quad Y \le y
\]

represents the event:

\[
\{X \le x \text{ and } Y \le y\}
\]

\subsection*{Definition}

The joint cumulative distribution function is defined as:

\[
F(x,y)=P(X\le x,Y\le y)
\]

This represents the probability that both random variables are less than or equal to specified values simultaneously.

\section{Joint CDF for Continuous Variables}

If $X$ and $Y$ have joint density $f(x,y)$ over region $[a,b]\times[c,d]$, then

\[
F(x,y)=\int_{c}^{y}\int_{a}^{x} f(u,v)\,du\,dv
\]

This means the CDF accumulates probability by integrating the density over the rectangle from the lower bounds to $(x,y)$.

\subsection*{Recovering the Joint PDF}

The joint density can be obtained from the CDF using partial derivatives:

\[
f(x,y)=\frac{\partial^2}{\partial x \partial y}F(x,y)
\]

Thus the CDF and PDF are mathematically related through differentiation and integration.

\section{Joint CDF for Discrete Variables}

If $X$ and $Y$ are discrete with joint PMF $p(x_i,y_j)$, then

\[
F(x,y)=\sum_{x_i \le x}\sum_{y_j \le y} p(x_i,y_j)
\]

Interpretation:

To compute $F(x,y)$:

\begin{itemize}
\item Select all rows with $x_i \le x$
\item Select all columns with $y_j \le y$
\item Add all probabilities in that region of the table.
\end{itemize}

\section{Properties of the Joint CDF}

The joint CDF $F(x,y)$ satisfies several important properties.

\subsection*{1. Non-decreasing property}

$F(x,y)$ is non-decreasing

If $x$ or $y$ increases, $F(x,y)$ either:

\begin{itemize}
\item stays the same, or
\item increases.
\end{itemize}

\subsection*{2. Lower-left boundary condition}

\[
F(x,y)=0
\]

at the lower-left corner of the joint range.

If the range extends to $(-\infty,-\infty)$:

\[
\lim_{(x,y)\to(-\infty,-\infty)}F(x,y)=0
\]

\subsection*{3. Upper-right boundary condition}

\[
F(x,y)=1
\]

at the upper-right corner.

If the range extends to $(\infty,\infty)$:

\[
\lim_{(x,y)\to(\infty,\infty)}F(x,y)=1
\]

These properties guarantee that the function behaves like a valid cumulative probability function.

\end{document}